\slajdid{09_2-s047}
\begin{frame}[label=09_2-s047]{Dimenzije tranzistora CMOS invertora - jedinični CMOS invertor}
\gusttekst\setlength{\abovedisplayskip}{3pt}\setlength{\belowdisplayskip}{3pt}\setlength{\abovedisplayshortskip}{2pt}\setlength{\belowdisplayshortskip}{2pt}Pojam jediničnog CMOS invertora će nam se pojavljivati i pojavljuje se kao reprezentativan primer za odgovarajuću tehnologiju. Pored ostalih parametara tehnologija određuje i sa kojim minimalnim dimenzijama raspolažemo, odnosno koliko su nam minimalne dužine i širine kanala. Oznaka tehnologije je istovremeno i ta minimalna dimenzija. Znači kada se kaže da je digitalni sistem, na primer procesor, urađen u 45nm tehnologiji, to znači da su tranzistori koji imaju minimalnu geometriju dužine kanala 45nm i širine kanala 45nm. Ono što je najčešće kao što smo i mi radili jeste da se parametri tranzistora podešavaju samo promenom dužine i širine kanala. U tom slučaju dužina i širina kanala su celobrojni multipl minimalnih dimenzija.\par\bigskip
(10µm–1971, \href{https://en.wikipedia.org/wiki/6_\%C2\%B5m_process}{{\color{DEplava}\underline{6µm}}}–1974, \href{https://en.wikipedia.org/wiki/3_\%C2\%B5m_process}{{\color{DEplava}\underline{3µm}}}–1977, \href{https://en.wikipedia.org/wiki/1.5_\%C2\%B5m_process}{{\color{DEplava}\underline{1.5µm}}}–1981, \href{https://en.wikipedia.org/wiki/1_\%C2\%B5m_process}{{\color{DEplava}\underline{1µm}}}–1984, \href{https://en.wikipedia.org/wiki/800_nm_process}{{\color{DEplava}\underline{800nm}}}–1987, \href{https://en.wikipedia.org/wiki/600_nm_process}{{\color{DEplava}\underline{600nm}}}–1990, \href{https://en.wikipedia.org/wiki/350_nm_process}{{\color{DEplava}\underline{350nm}}}–1993, \href{https://en.wikipedia.org/wiki/250_nm_process}{{\color{DEplava}\underline{250nm}}}–1996, \href{https://en.wikipedia.org/wiki/180_nm_process}{{\color{DEplava}\underline{180nm}}}–1999, \href{https://en.wikipedia.org/wiki/130_nm_process}{{\color{DEplava}\underline{130nm}}}–2001, \href{https://en.wikipedia.org/wiki/90_nm_process}{{\color{DEplava}\underline{90nm}}}–2003, \href{https://en.wikipedia.org/wiki/65_nm_process}{{\color{DEplava}\underline{65nm}}}–2005, 45nm–2007, \href{https://en.wikipedia.org/wiki/32_nm_process}{{\color{DEplava}\underline{32nm}}}–2009, \href{https://en.wikipedia.org/wiki/22_nm_process}{{\color{DEplava}\underline{22nm}}}–2012, \href{https://en.wikipedia.org/wiki/14_nm_process}{{\color{DEplava}\underline{14nm}}}–2014, \href{https://en.wikipedia.org/wiki/10_nm_process}{{\color{DEplava}\underline{10nm}}}–2016, \href{https://en.wikipedia.org/wiki/7_nm_process}{{\color{DEplava}\underline{7nm}}}–2018, \href{https://en.wikipedia.org/wiki/5_nm_process}{{\color{DEplava}\underline{5nm}}}–\textasciitilde{}2020)
\end{frame}
